% !TEX root = ../report.tex
\structchapter{ВВЕДЕНИЕ}
Линейный и бинарный поиск решают одну задачу, но различаются числом сравнений, дополнительной памятью и требованиями к массиву. В работе сравниваются четыре реализации: линейная, две нерекурсивные бинарные и рекурсивная бинарная.

Цель работы --- выполнить сравнительный анализ ресурсной эффективности алгоритмов поиска в массиве.

Для достижения цели необходимо решить следующие задачи:
\begin{enumerate}
\item описать задачу поиска значения в массиве;
\item разработать линейный, два нерекурсивных бинарных и рекурсивный бинарный алгоритмы;
\item представить разработанные алгоритмы схемами;
\item сформулировать требования к программному обеспечению;
\item выбрать язык программирования и инструменты сборки;
\item реализовать алгоритмы и программу для их сравнения;
\item провести тестирование реализаций;
\item оценить дополнительную память каждого алгоритма;
\item оценить число сравнений в лучшем и худшем случаях;
\item представить результаты сравнительного анализа в таблицах;
\item сформулировать выводы о числе сравнений, затратах памяти и ограничениях алгоритмов.
\end{enumerate}

